\documentclass[10pt,a4paper]{amsart}
\usepackage{amsmath,amssymb,amsthm,graphicx}
% For commutative diagrams we use Xy-pic
\usepackage[all]{xy}

% We use 2cell for 2-commutative diagrams.
\xyoption{2cell}
% Theorem environments.
%
\theoremstyle{plain}
\newtheorem{theorem}[subsection]{Theorem}
\newtheorem{proposition}[subsection]{Proposition}
\newtheorem{lemma}[subsection]{Lemma}

\theoremstyle{definition}
\newtheorem{definition}[subsection]{Definition}
\newtheorem{example}[subsection]{Example}
\newtheorem{exercise}[subsection]{Exercise}
\newtheorem{situation}[subsection]{Situation}

\theoremstyle{remark}
\newtheorem{remark}[subsection]{Remark}
\newtheorem{remarks}[subsection]{Remarks}

\numberwithin{equation}{subsection}

% Macros
%
\def\lim{\mathop{\mathrm{lim}}\nolimits}
\def\colim{\mathop{\mathrm{colim}}\nolimits}
\def\Spec{\mathop{\mathrm{Spec}}}
\def\Hom{\mathop{\mathrm{Hom}}\nolimits}
\def\Ext{\mathop{\mathrm{Ext}}\nolimits}
\def\SheafHom{\mathop{\mathcal{H}\!\mathit{om}}\nolimits}
\def\SheafExt{\mathop{\mathcal{E}\!\mathit{xt}}\nolimits}
\def\Sch{\mathit{Sch}}
\def\Mor{\mathop{\mathrm{Mor}}\nolimits}
\def\Ob{\mathop{\mathrm{Ob}}\nolimits}
\def\Sh{\mathop{\mathit{Sh}}\nolimits}
\def\NL{\mathop{N\!L}\nolimits}
\def\CH{\mathop{\mathrm{CH}}\nolimits}
\def\proetale{{pro\text{-}\acute{e}tale}}
\def\etale{{\acute{e}tale}}
\def\QCoh{\mathit{QCoh}}
\def\Ker{\mathop{\mathrm{Ker}}}
\def\Im{\mathop{\mathrm{Im}}}
\def\Coker{\mathop{\mathrm{Coker}}}
\def\Coim{\mathop{\mathrm{Coim}}}

% Boxtimes
%
\DeclareMathSymbol{\boxtimes}{\mathbin}{AMSa}{"02}

%
% Macros for moduli stacks/spaces
%
\def\QCohstack{\mathcal{QC}\!\mathit{oh}}
\def\Cohstack{\mathcal{C}\!\mathit{oh}}
\def\Spacesstack{\mathcal{S}\!\mathit{paces}}
\def\Quotfunctor{\mathrm{Quot}}
\def\Hilbfunctor{\mathrm{Hilb}}
\def\Curvesstack{\mathcal{C}\!\mathit{urves}}
\def\Polarizedstack{\mathcal{P}\!\mathit{olarized}}
\def\Complexesstack{\mathcal{C}\!\mathit{omplexes}}
% \Pic is the operator that assigns to X its picard group, usage \Pic(X)
% \Picardstack_{X/B} denotes the Picard stack of X over B
% \Picardfunctor_{X/B} denotes the Picard functor of X over B
\def\Pic{\mathop{\mathrm{Pic}}\nolimits}
\def\Picardstack{\mathcal{P}\!\mathit{ic}}
\def\Picardfunctor{\mathrm{Pic}}
\def\Deformationcategory{\mathcal{D}\!\mathit{ef}}
\usepackage{cleveref}
\crefname{originallemma}{original lemma}{original lemmas}
\Crefname{originallemma}{Original lemma}{Original lemmas}
\makeatletter
\newlabel{cohomology-lemma-irreducible-constant-cohomology-zero}{{02UW}{source}{}{sourcepointer-1}{}}
\newlabel{cohomology-lemma-irreducible-constant-cohomology-zero@cref}{{[originallemma][0][]02UW}{[1][0][]source}{}{}{}}
\makeatother
\crefname{originalreference}{original reference}{original references}
\Crefname{originalreference}{Original reference}{Original references}
\makeatletter
\newlabel{lemma-higher-direct-images}{{03Q8}{source}{}{sourcepointer-2}{}}
\newlabel{lemma-higher-direct-images@cref}{{[originalreference][0][]03Q8}{[1][0][]source}{}{}{}}
\makeatother
\crefname{originalreference}{original reference}{original references}
\Crefname{originalreference}{Original reference}{Original references}
\makeatletter
\newlabel{lemma-local-rings-strictly-henselian}{{09AX}{source}{}{sourcepointer-3}{}}
\newlabel{lemma-local-rings-strictly-henselian@cref}{{[originalreference][0][]09AX}{[1][0][]source}{}{}{}}
\makeatother
\crefname{originalreference}{original reference}{original references}
\Crefname{originalreference}{Original reference}{Original references}
\makeatletter
\newlabel{lemma-normal-scheme-with-alg-closed-function-field}{{09Z9}{source}{}{sourcepointer-4}{}}
\newlabel{lemma-normal-scheme-with-alg-closed-function-field@cref}{{[originalreference][0][]09Z9}{[1][0][]source}{}{}{}}
\makeatother
\title{Higher direct images along a morphism from an integral normal scheme with separably closed function field vanish}
\author{The Stacks Project Authors}
\date{stacks-project:etale-cohomology.tex:local-snapshot-20261011; GNU FDL 1.2 or later}
\begin{document}
\maketitle
\noindent\textit{Editorial note.} All delivered mathematical text is the author's own, copied from the pinned Stacks Project chapter file etale-cohomology.tex; no statement or proof step was rewritten, reordered or removed.\par\medskip
\noindent\textit{Editorial note.} The counted proof is the author's complete argument: the proof is written out in full: the presheaf description of the higher direct images, the decomposition of the base change into open subschemes, and the two inputs giving the vanishing.\par\medskip
\noindent\textit{Editorial note.} The label lemma-Rf-star-zero-normal-with-alg-closed-function-field names the counted statement, and the proof environment immediately following it in the source is delivered verbatim as its proof.\par\medskip
\noindent\textit{Editorial note.} The counted statement is carried by the author's own lemma environment in the chapter source.\par\medskip
\noindent\textit{Editorial note.} The statement and proof are copied from the pinned chapter file; the four cited inputs are kept as external pointers with their canonical Stacks tags.\par\medskip
\noindent\textit{Allowed standard prerequisites (editorial).}\begin{itemize}
\item Higher direct images as sheafified cohomology of the base change. Same chapter, the lemma describing R\textasciicircum{}q f\_*underline M as the sheaf associated to V |-> H\textasciicircum{}q\_etale(V x\_Y X, M). Supplies the presheaf description used at the start.
\item Structure of separated etale morphisms over such a scheme. Same chapter, the lemma stating that a separated etale morphism over an integral normal scheme with separably closed function field is a disjoint union of open immersions. Supplies the decomposition V x\_Y X = coproduct of open subschemes U\_i.
\item Equality of etale and Zariski cohomology and vanishing on an irreducible scheme. Same chapter, the lemma identifying etale and Zariski cohomology when all local rings are strictly henselian, together with Cohomology, the vanishing of the cohomology of a constant sheaf on an irreducible space. Supplies the vanishing of H\textasciicircum{}q\_etale(U\_i, M) for q > 0.
\end{itemize}\medskip
\ifdefined\hypertarget\hypertarget{sourcepointer-1}{}\fi
\noindent\textit{Source reference.} Original lemma 02UW: sources/stacks-chow/tags.txt, source lines 3747--3747. Prerequisite pointer invoked at proof line(s) 14386; it names a standard fact of a sibling chapter that the counted argument uses, and every step of the counted proof built on it is retained verbatim in the delivered body.\par\medskip
\ifdefined\hypertarget\hypertarget{sourcepointer-2}{}\fi
\noindent\textit{Source reference.} Original reference 03Q8: sources/stacks-etale-cohomology/pinned-original-source.tex, source lines 7510--7510. Prerequisite pointer invoked at proof line(s) 14378; it names a standard fact of the same chapter that the counted argument uses, and every step of the counted proof built on it is retained verbatim in the delivered body.\par\medskip
\ifdefined\hypertarget\hypertarget{sourcepointer-3}{}\fi
\noindent\textit{Source reference.} Original reference 09AX: sources/stacks-etale-cohomology/pinned-original-source.tex, source lines 14205--14205. Prerequisite pointer invoked at proof line(s) 14383; it names a standard fact of the same chapter that the counted argument uses, and every step of the counted proof built on it is retained verbatim in the delivered body.\par\medskip
\ifdefined\hypertarget\hypertarget{sourcepointer-4}{}\fi
\noindent\textit{Source reference.} Original reference 09Z9: sources/stacks-etale-cohomology/pinned-original-source.tex, source lines 14312--14312. Prerequisite pointer invoked at proof line(s) 14382; it names a standard fact of the same chapter that the counted argument uses, and every step of the counted proof built on it is retained verbatim in the delivered body.\par\medskip
\begin{lemma}
\label{lemma-Rf-star-zero-normal-with-alg-closed-function-field}
Let $f : X \to Y$ be a morphism of schemes where $X$ is an integral
normal scheme with separably closed function field. Then
$R^qf_*\underline{M} = 0$ for $q > 0$ and any abelian group $M$.
\end{lemma}

\begin{proof}
Recall that $R^qf_*\underline{M}$ is the sheaf associated
to the presheaf $V \mapsto H^q_\etale(V \times_Y X, M)$ on $Y_\etale$, see
Lemma \ref{lemma-higher-direct-images}.
If $V$ is affine, then $V \times_Y X \to X$ is separated and \'etale.
Hence $V \times_Y X = \coprod U_i$ is a disjoint union of open
subschemes $U_i$ of $X$, see
Lemma \ref{lemma-normal-scheme-with-alg-closed-function-field}.
By Lemma \ref{lemma-local-rings-strictly-henselian}
we see that $H^q_\etale(U_i, M)$ is equal to
$H^q_{Zar}(U_i, M)$. This vanishes by
Cohomology, Lemma \ref{cohomology-lemma-irreducible-constant-cohomology-zero}.
\end{proof}

\end{document}
